%
%	z.B. Text des ersten Autors
%

\section{Einführung}
Eine Untergliederung des Dokuments in Abschnitte und Unterabschnitte
geschieht mit den Kommandos \verb|\section{...}| und \verb|\subsection{...}|,
in Absätze mit \verb|\par| oder einfach einer Leerzeile im \LaTeX-Quelltext.

\paragraph{Betitelter Absatz.}
Ein sehr kleiner Abschnitt (meistens nur ein Absatz)
mit Titeltext, der nicht vom Textrumpf abgesetzt ist, wird durch das Kommando
\verb|\paragraph{titel}| erzeugt.

Bitte seien Sie mit dem Kommando \verb?\verb|text|? vorsichtig! Die Anweisung bewirkt eine exakte Darstellung des Textes darin.

%
%	z.B. Text des zweiten Autors
%

\section{Weitere Kommandos}

\subsection{Mathematische Formeln}\label{sec:formeln}

Mathematische Formeln werden mittels \verb|\(...\)|
in den Fließtext eingebaut
--- zum Beispiel \( E=mc^2 \) und:
sei \(V\) ein Vektorraum über \(\mathbb{R}\)
und \(\mathcal{M}\) eine Indexmenge
---
oder aber mittels \verb|\[...\]|
abgesetzt und zentriert dargestellt:
\[
    \pmb{x} = \sqrt[3]{\frac{a^2-b^2}{a^2+b^2}}
    ~~~~~ \text{versus} ~~~~~
    \boldsymbol{x} = \sqrt[3]{\frac{a^2-b^2}{a^2+b^2}}
\]


\subsection{Silbentrennung}\label{sec:silbentrennung}
Vertrauen Sie bitte nie einer automatischen Silbentrennung (auch nicht der von Microsoft Word \& Co.). In folgendem Test-Absatz ist das Wort "`Spracherkennung"' falsch getrennt.

Testzeile Testzeile Testzeile Testzeile Testzeile Testzeile Testzeile Testzeile Te Spracherkennung.

Sie können LaTeX die richtige Trennung mit \verb|\hyphenation{...}| mitteilen. Man tut das üblicherweise noch vor \verb|\begin{document}|.

\hyphenation{Sprach-er-ken-nung}	% hier jetzt zu Demonstrationszwecken im Text
Testzeile Testzeile Testzeile Testzeile Testzeile Testzeile Testzeile Testzeile Te Spracherkennung.


\subsection{Literaturzitate}
\label{sec:literaturzitate}

Diese Vorlage basiert auf einer Latex Vorlage von Claudius Stern\indircite{Stern2018b} mit meinen eigenen Anpassungen.
Konkret wurden extra Zitierbefehle angelegt, um angepasst an die Anforderungen aus dem Kurs Wissenschaftliches Arbeiten zitieren zu können.

Das Literaturverwaltungsprogramm JabRef\indircite{JabRef2014} ist für viele Plattformen verfügbar und unterstützt bei der Literaturrecherche. Es ist prädestiniert dazu, mit {\LaTeX} in Kombination mit {Bib\TeX} zusammenzuarbeiten.

So kann in dieser Vorlage mit \verb|\dircite{Bibtexkey}| eine Quelle direkt Zitiert werden.
Optional kann auch ein Seitenhinweis mit \verb|\dircite[S. 42]{Buch}| wie folgt beigefügt werden.
Indirekte Zitate können analog mit \verb|\indircite[page info]{Bibtexkey}| erstellt werden.

\subsection{Einbinden von Bildern}

Sie können mit der Anweisung \verb|\includegraphics{datei}| eine Grafikdatei einbinden. Diese kann im PDF-, JPG- oder PNG-Format vorliegen. Grafiken werden üblicherweise in die float-Umgebung \verb|figure| gekapselt. Jede eingebundene Grafik muss vom Text aus referenziert werden. Die textuelle Referenz hat \textit{vor} der Grafik zu erfolgen.

\begin{figure}[tb]
    \centering
    \includegraphics{images/FOM.jpg}
    \dircitegraph{FOM}
    \caption{Logo der FOM}
    \label{fig:fomlogo}
\end{figure}

Auch andere Grafikformate werden unterstützt.
Grafiken finden sich nicht notwendigerweise direkt am Einfüge-Ort (siehe \autoref{fig:fomlogo}).
Das Textsatzsystem richtet es so ein, dass es gut aussieht.
Zitate von Graphiken sind möglich über \verb|\dircitegraph[]{}| und \verb|\indircitegraph[]{}|.

\subsection{Querverweise}\label{sec:querverweise}
Querverweise können sie mit dem \verb|\label{key}| befehl direkt nach einer section Anweisung oder als Teil einer Graphik erstellen.
Auf diese können sie dann anschließend \verb|\autoref{key}|, oder wenn sie nur die Nummer brauchen mit \verb|\ref{key}| verweisen.

Näheres zu Formeln können sie Beispielsweise in \autoref{sec:formeln} lesen, während in \ref{sec:silbentrennung} Silbentrennung erklärt wird.

\subsection{Abkürzungen}
Abkürzung können nach einem Eintrag in der Datei \verb|abbreviation.tex| mit den Befehlen \verb|\ac{XYZ}| (für einen kleinen ersten Buchstaben) und \verb|\Ac{XYZ}| (für einen großen ersten Buchstaben) in den Text eingefügt werden.

Dieses Template ist für Seminararbeiten an der \ac{FOM} gedacht.
Es wurde für einen \ac{IT} Studiengang der \ac{FOM} angepasst.

\subsection{Erstellen von Tabellen}

Das Volk hat gesprochen. Siehe \autoref{tab:tabular}. Auch hier kommt ein float zum Einsatz, jedoch mit dem Positionierungs-Hinweis \verb|[h]|. Jedes float sollte mit einem Label versehen werden, und es sollte im Text darauf verwiesen werden, da sich die Position ändern kann.

\begin{table}[h]
    \centering
    \begin{tabular} {|llc||r|}
        \hline
        Name       & Rang    & Fraktion & Stimmenanteil \\
        \hline
        Mobutu     & General & CDU      & 57\%          \\
        Tsvangirai & Oberst  & CSU      & 63\%          \\
        \hline
    \end{tabular}
    \caption {Bundestagswahl in Simbabwe}
    \label{tab:tabular}
\end{table}

\autoref{tab:booktab} verwendet keine vertikalen Linien und entspricht dem üblicherweise in Büchern verwenden Stil.
Derartige Tabellen sind deutlich ansehnlicher.

\begin{table}[h]
    \centering
    \begin{tabular} {llcr}
        \toprule
        Name       & Rang    & Fraktion & Stimmenanteil \\
        \midrule
        Mobutu     & General & CDU      & 57\%          \\
        Tsvangirai & Oberst  & CSU      & 63\%          \\
        \bottomrule
    \end{tabular}
    \caption {Bundestagswahl in Simbabwe}
    \label{tab:booktab}
\end{table}

\subsection{Listen und Aufzählungen}

Listen und Aufzählungen werden in einer Umgebung angelegt (umschlossen von \verb|begin| und \verb|end|.).
\verb|\begin{itemize}| leitet eine Liste ein und \verb|\begin{enumerate}| eine Aufzählung.
Die Einträge werden jeweils mit \verb|\item| begonnen.
Folgend zwei Beispiele.

Itemize:
\begin{itemize}
    \item Test
    \item Test
    \item Test
\end{itemize}

Enumerate:
\begin{enumerate}
    \item Test
    \item Test
    \item Test
\end{enumerate}
